The second contribution concerns the annealed record-moment free energy and its replica phase diagram, not a universal closure of normalized trajectories. Exact finite-size spectra are supplemented by thermodynamic-limit uniform-MPS calculations~\cite{ZaunerStauber2018,Fishman2018}. At $n=3$ these give $p_c^{(3)}=0.3049(2)$ and simultaneous loss of sign and standard order; weak first order is the preferred reading, with a continuous alternative still open. The earlier $q=n$-Potts assignments at $n=3,4$ are withdrawn. Their failure motivates a structural reinterpretation: the globally aligned permutation manifold carries the regular $S_n$-bimodule and has $n!$ labels, connecting to the permutation-Potts description of Ref.~\cite{Vasseur2019}. Neither the number of labels nor the dimension of one irrep establishes a universality class. At $n=4$ the transition location is only approximately constrained, around $0.38$--$0.39$; at $n=5$ the $L=4,6,8$ closing diagnostics are compatible with stronger first-order behavior, not a completed asymptotic test.

The third contribution is logically independent of those class assignments. Born normalization separates conditional trajectories from symmetric replica summaries: an information-theoretic no-go has an explicit bivariate recovery, and conditional-trajectory and conditional-purity decision are PostBQP-complete for universal-circuit descriptions under the promises of Appendix~\ref{app:complexity}. The record theory supplies quenched large-deviation and multifractal identities, a no-freezing statement with explicit analytic-cocycle and nondegeneracy hypotheses, the exact replica interpolation identity, and numerical quenched Haar record SCGFs. These are results in their own right, not merely consequences of a proposed phase diagram. Clifford records provide exact dyadic identities and a second numerical route: the Haar $n=3$ annealed moment equals a positive $4^{-N_{\rm rand}}$ disorder tilt of stabilizer circuits. We implement that route and state the finite-population and genealogy limitations of its present pilot.

\emph{Version correction (16 October 2026).} This author-review revision withdraws the v32 claims of three-state Potts at $n=3$, four-state Potts with logarithmic corrections at $n=4$, and the $q=n$ class-identification argument. The finite-size numerical values are retained, with their interpretation corrected. The exact $n=2$ Ising result is credited to Bao, Choi and Altman~\cite{Bao2020}; the new work here concerns operator contracts, rank and spectral certificates, the extended computations and record theory. Section~\ref{sec:n3annealed} includes a retrospective rather than hiding the failed inference. New numerical evidence is not promoted to a theorem.

A delimitation is essential. The integer-$n$ operators compute annealed moments, not the Born-weighted entanglement transition. Appendix~\ref{app:numerics} reports the author's separate Clifford finite-size-scaling benchmark $p_c=0.1597(8)$, $\nu=1.24(7)$; the Haar benchmark of Ref.~\cite{Zabalo2020} is $p_c\simeq0.17(1)$ for its circuit conventions. Fits from a few annealed rungs toward $n=1$ are suggestive extrapolations only and distinguish neither benchmark. The information-theoretic no-go remains in force: reorganizing the paper around one annealed function does not make that function sufficient for normalized entanglement.

The logical structure is as follows. Sections~\ref{sec:compression}--\ref{sec:ising} establish the exact two-replica mechanics and recover the Ising solution; Sections~\ref{sec:general-n}--\ref{sec:n4n5} give general-$n$ structure and the corrected numerical ladder. Section~\ref{sec:born} states the Born-versus-symmetric distinction and complexity results. Section~\ref{sec:records} develops record statistics, interpolation, the Clifford identities and the new tilted-SMC pilot. Section~\ref{sec:aux} delimits the continuum program. Sections~\ref{sec:interpretation} and~\ref{sec:conclusion} distinguish operational consequences, interpretations and open tests.
